% summary.tex -- a guide to the Lean formalization in percolation/ of anthropics/formal-math.
% Convention for this repository: one summary.tex per project, article class, sections
%   1 Introduction and history / 2 Main results (each with its Lean name) / 3 Proof sketch /
%   4 Guide to the formalization / 5 Limitations; references.bib = the project's references.bib
%   plus the entries cited only here.  The warrant for every theorem is the Lean development;
%   this text is a reader's guide to it and has not been independently refereed.
\documentclass[11pt]{article}
\usepackage[T1]{fontenc}
\usepackage{lmodern}
\usepackage[margin=0.86in]{geometry}
\usepackage{amsmath,amssymb,amsthm,mathtools}
\usepackage{enumitem,booktabs,array,xcolor}
\usepackage[obeyspaces,spaces,hyphens]{url}
\usepackage[numbers,sort&compress]{natbib}
\usepackage[colorlinks=true,linkcolor=blue!40!black,citecolor=green!35!black,urlcolor=blue!55!black,pdftitle={theta(p_c)=0 for Bernoulli bond percolation on Z^d, d>=2: a guide to the Lean formalization},pdfauthor={Justin Leder}]{hyperref}
\usepackage[capitalise,noabbrev]{cleveref}
\usepackage[small,compact]{titlesec}

\newtheorem{theorem}{Theorem}[section]
\newtheorem{corollary}[theorem]{Corollary}
\newtheorem{lemma}[theorem]{Lemma}
\newtheorem{proposition}[theorem]{Proposition}
\theoremstyle{definition}
\newtheorem{definition}[theorem]{Definition}
\theoremstyle{remark}
\newtheorem{remark}[theorem]{Remark}

% notation
\newcommand{\PP}{\mathbb{P}}
\newcommand{\E}{\mathbb{E}}
\DeclareMathOperator{\Cov}{Cov}
\newcommand{\Zd}{\mathbb{Z}^d}
\newcommand{\ZZ}{\mathbb{Z}}
\newcommand{\RR}{\mathbb{R}}
\newcommand{\pc}{p_{\mathrm c}}
\newcommand{\con}{\leftrightarrow}
\newcommand{\ncon}{\nleftrightarrow}
\newcommand{\one}{\mathbf 1}
\newcommand{\eC}{\mathcal{C}}
\newcommand{\Sur}{\operatorname{Sur}}
\newcommand{\Marg}{\operatorname{Marg}}
\newcommand{\slf}{\operatorname{sl}}
\newcommand{\HH}{\mathbb{H}}
\newcommand{\CSH}{\mathrm{CSH}}
\newcommand{\GEN}{\mathrm{GEN}}
% Lean identifiers and file paths (line-breakable typewriter)
\DeclareUrlCommand\lean{\def\UrlFont{\ttfamily\upshape}}

% allow \lean identifiers to break also after underscores and capitals-free long runs
\makeatletter
\g@addto@macro\UrlBreaks{\do\_\do\a\do\b\do\c\do\d\do\e\do\f\do\g\do\h\do\i\do\j\do\k\do\l\do\m\do\n\do\o\do\p\do\q\do\r\do\s\do\t\do\u\do\v\do\w\do\x\do\y\do\z}
\makeatother
\emergencystretch=1.5em
\makeatletter\g@addto@macro\normalsize{\setlength\abovedisplayskip{6.5pt plus 2pt minus 3pt}\setlength\belowdisplayskip{6.5pt plus 2pt minus 3pt}%
\setlength\abovedisplayshortskip{1pt plus 2pt}\setlength\belowdisplayshortskip{4pt plus 2pt minus 2pt}}\makeatother
\setlist{itemsep=0pt,topsep=2pt,parsep=0pt,partopsep=0pt}
\setlength{\textfloatsep}{12pt plus 2pt minus 3pt}

\title{\vspace{-2.2em}\texorpdfstring{$\theta(\pc)=0$}{theta(pc)=0} for Bernoulli bond percolation on \texorpdfstring{$\Zd$}{Zd}\\ in all dimensions \texorpdfstring{$d\ge 2$}{d>=2}:\\[3pt]
\large a guide to the Lean formalization}
\author{Justin Leder}
\date{\vspace{-0.3em}August 2026\vspace{-0.8em}}

\begin{document}
\maketitle

\begingroup\renewcommand\thefootnote{}\footnotetext{%
\emph{Provenance and review status.}
The Lean sources of this development---definitions, statements and proofs, together with the statement file \texttt{Challenge.lean}, its
proved twin \texttt{Solution.lean} and the metadata---were written by an AI system (Anthropic's Claude models) working under my direction;
no human wrote or edited the Lean code, and the first drafts of this guide were produced in the same way.  Correctness rests on mechanical
checking: the Lean~4 kernel accepts every file with no \texttt{sorry} outside the two deliberate placeholders of \texttt{Challenge.lean}, no
added axioms and no unsafe code, and the main theorems depend only on the standard axioms \texttt{propext}, \texttt{Classical.choice},
\texttt{Quot.sound}; the compared theorems are additionally replayed in the independent \texttt{nanoda} kernel by the Lean comparator (see
\texttt{AUDIT.md}).  Neither the development nor this guide has yet been refereed by human mathematicians or by anyone independent of me; the
only review so far was carried out by AI systems.  Mechanical checking does not cover whether the formal statements express the intended
mathematics: readers should satisfy themselves that \texttt{Challenge.lean} states the theorem of \cref{s:results} (box ``Statements to
audit'').}\endgroup

\begin{abstract}
\noindent
This is a reader's guide to a Lean~4/Mathlib formalization proving that the percolation probability of nearest-neighbour Bernoulli bond
percolation on $\Zd$ vanishes at the critical point, $\theta(\pc)=0$, for every $d\ge2$; the cases $3\le d\le10$, in particular $\ZZ^3$, were
open.  The route is the reduction of Kozma and Nitzan (2024), who conjectured a family of
``gluing'' inequalities for percolation on arbitrary finite weighted graphs and proved that the weakest of them, their Conjecture~3,
implies $\theta(\pc)=0$ on $\Zd$ for all $d\ge2$.  The development proves Conjecture~3 through a stronger additive gluing
inequality---if $\PP(a\con b)\ge1-t$ for every $a\in A$ then $\PP(o\con b)\ge\PP(o\con A)-t$---which is in turn derived from a new
family of conditioned covariance inequalities for increasing functions of a single open cluster, indexed by finite lists of auxiliary
vertices and proved by induction on the list; its first member is the Harris inequality.  Kozma--Nitzan's Theorem~6 and every
classical input are re-proved inside the library, so the final statement has no hypothesis other than $d\ge2$.  The warrant for every claim
is the Lean development, not this text.
\end{abstract}

%=====================================================================================================
\section{Introduction}\label{s:intro}

\subsection{The model and the question}\label{ss:model}
Bernoulli bond percolation on $\Zd$, introduced by Broadbent and Hammersley \citep{BroadbentHammersley1957}, declares every
nearest-neighbour edge of $\Zd$ \emph{open} with probability $p\in[0,1]$, independently; $\PP_p$ is the product measure and $C_0$ the
set of vertices joined to the origin by open paths.  The \emph{percolation probability} is $\theta(p):=\PP_p(|C_0|=\infty)$ and the
\emph{critical probability} $\pc=\pc(\Zd):=\inf\{p:\theta(p)>0\}$ satisfies $0<\pc<1$ for $d\ge2$
\citep[Theorem~(1.10)]{GrimmettPercolation1999}.  Since $\theta$ is non-decreasing, vanishes on $[0,\pc)$, and is continuous on
$(\pc,1]$ \citep{VandenBergKeane1984}, \citep[Theorem~(8.8)]{GrimmettPercolation1999}---the last fact resting on the uniqueness of the
infinite cluster \citep{AizenmanKestenNewmanCMP1987,BurtonKeane1989}---$\theta$ is continuous on $[0,1]$ if and only if
\begin{equation}\label{eq:continuity}
 \theta(\pc)=0,
\end{equation}
i.e.\ if and only if there is almost surely no infinite cluster at the critical point \citep[\S8.3]{GrimmettPercolation1999}.
That \eqref{eq:continuity} holds for every $d\ge2$ is conjectured---``open since at least the 80s'' \citep[p.~1]{KozmaNitzan2024}---and
recorded as open in \citep[pp.~14, 202--203]{GrimmettPercolation1999}, \citep[Conjecture~1]{DuminilCopin2018ICM}.  This guide describes a formal proof of \eqref{eq:continuity} for all $d\ge2$.

\subsection{What was known}\label{ss:history}
\emph{Two dimensions.}  Harris \citep{HarrisPCPS1960} proved $\theta(\tfrac12)=0$ on $\ZZ^2$ and Kesten \citep{KestenCMP1980} proved
$\pc(\ZZ^2)=\tfrac12$, building on Russo \citep{Russo1978} and Seymour--Welsh \citep{SeymourWelsh1978}; together these give
\eqref{eq:continuity} for $d=2$.  Harris' paper also contains the correlation inequality that bears his name (in general form
Fortuin--Kasteleyn--Ginibre \citep{FortuinKasteleynGinibreCMP1971}), which with the van den Berg--Kesten inequality \citep{VandenBergKesten1985}
is the basic tool of the subject.  \emph{Sharpness}---exponential decay of the cluster radius for $p<\pc$, due to Menshikov
\citep{Menshikov1986} and Aizenman--Barsky \citep{AizenmanBarsky1987}, with short proofs by Duminil-Copin--Tassion
\citep{DuminilCopinTassionCMP2016,DuminilCopinTassionEM2016}---does not decide \eqref{eq:continuity}.
\emph{High dimensions.}  Aizenman--Newman \citep{AizenmanNewman1984} introduced the triangle condition and Barsky--Aizenman
\citep{BarskyAizenman1991} showed that it implies \eqref{eq:continuity}; Hara--Slade \citep{HaraSlade1990} verified it by the lace expansion in
high dimension ($d\ge19$ \citep[Thm.~2.7]{HaraSlade1994}), and Fitzner--van der Hofstad \citep{FitznerVanDerHofstad2017} reached $d\ge11$.
Thus \eqref{eq:continuity} was known for $d=2$ and $d\ge11$ and open for $3\le d\le10$.
\emph{Restricted geometries.}  Barsky--Grimmett--Newman \citep{BarskyGrimmettNewman1991} proved by a block argument at fixed $p$ that a
half-space does not percolate at its critical point, which by Grimmett--Marstrand \citep{GrimmettMarstrand1990} equals $\pc(\Zd)$, so
$\theta_{\HH}(\pc)=0$ \citep[Thms.~(7.2), (7.35)]{GrimmettPercolation1999}; Duminil-Copin--Sidoravicius--Tassion
\citep{DuminilCopinSidoraviciusTassion2016} proved that slabs $\ZZ^2\times\{0,\dots,k\}$ do not percolate at their own critical points.  As
Kozma and Nitzan stress \citep[\S1]{KozmaNitzan2024}, in the block arguments the step from a finite-size criterion back to percolation needs an
extra feature (the half-space, or an increase of $p$); that step in $\Zd$ itself, at fixed $p$, was missing.

\subsection{The Kozma--Nitzan reduction}\label{ss:KN}
Kozma and Nitzan \citep{KozmaNitzan2024} proposed to supply the missing step by an inequality for percolation on \emph{finite} graphs
with arbitrary edge probabilities.  For such a graph, vertices $o,b$ and a vertex set $A$, with $\{o\con A\}$ the event that $o$ is
joined to some vertex of $A$, their Conjecture~1 \citep[p.~3]{KozmaNitzan2024} is
\begin{equation}\label{eq:KN1}
 \PP(o\con b)\ \ge\ \PP(o\con A)\cdot\min_{a\in A}\PP(a\con b),
\end{equation}
of which they write: ``We were not able to prove or disprove this conjecture.  Our belief that it holds is based on some (admittedly
restricted) numerical evidence and on some simple cases where we were able to prove it.''  They prove a stronger ``pre-FKG'' form for $|A|=2$, for
several configurations with $|A|=3$, and for $o$ close to $A$ \citep[Theorems~1--5]{KozmaNitzan2024}; they formulate variants (Conjecture~2;
Conjecture~4 for monotone functions of the cluster \citep[p.~32]{KozmaNitzan2024}); and they isolate the weakest statement of the family,
Conjecture~3 \citep[p.~15]{KozmaNitzan2024}:
\emph{for every $\varepsilon>0$ there is $\delta>0$ such that, for any finite weighted graph and any $A,o,b$, $\PP(o\con A)>1-\delta$ and
$\PP(a\con b)>1-\delta$ for all $a\in A$ imply $\PP(o\con b)>1-\varepsilon$.}  The content is that $\delta$ does not depend on $|A|$ (for bounded
$|A|$ the union bound gives $\delta=\varepsilon/(|A|+1)$).  Their main theorem is the reduction
\citep[Theorem~6]{KozmaNitzan2024}: \emph{if Conjecture~3 holds, then $\theta(\pc)=0$ on $\Zd$ for every $d\ge2$}, proved by a
one-step renormalisation at fixed $p$ (\cref{ss:thm6}).  The paper is an arXiv preprint (v1, January 2024); Gladkov
\citep[\S1]{Gladkov2024} reports the conjectured inequality as open, and we are not aware of a proof in print of any of Conjectures~1--4
beyond the cases above.

\subsection{What the formalization adds, and what it does not}\label{ss:adds}
The development proves Conjecture~3 (\cref{cor:conj3}) through the \emph{additive gluing inequality} (\cref{thm:AG}): if
$\PP(a\con b)\ge1-t$ for all $a\in A$ then $\PP(o\con b)\ge\PP(o\con A)-t$---weaker than \eqref{eq:KN1}, since
$\PP(o\con A)(1-t)\ge\PP(o\con A)-t$, and stronger than Conjecture~3 (take $t=\delta=\varepsilon/2$).  Additive gluing is derived from a family
of conditioned covariance inequalities for increasing functions of a single open cluster, the \emph{conditioned slack hierarchy}
(\cref{thm:CSH}), the one new object; its level-zero, unconditioned member is the Harris inequality.  Kozma--Nitzan's Theorem~6 with everything
it invokes, and every correlation inequality used, are re-proved inside the library, so the formal statement carries no hypothesis other than
$2\le d$.  Our contribution is the finite-graph inequality; the reduction, and the insight that an inequality uniform in $|A|$ is what the
renormalisation approach requires, are Kozma and Nitzan's.

\emph{Not claimed.}  Conjectures~1 \eqref{eq:KN1}, 2 and 4 are neither proved nor stated in the release (the three-relay case of
\eqref{eq:KN1} is proved, \cref{rem:mult}).  Nothing is proved about the rate at which $\theta(p)\to0$, critical exponents or the value of
$\pc$; nothing new about site percolation, other lattices, or slabs (the slab theorem of \citep{DuminilCopinSidoraviciusTassion2016} is
formalised as a literature input only).  Continuity of $\theta$ on $[0,1]$, classically equivalent to \eqref{eq:continuity}, is not part of
the formal statement.  No lace expansion, triangle condition or numerics are used.

\subsection{How to read this guide, and what to verify}\label{ss:howto}
The repository\footnote{\url{https://github.com/anthropics/formal-math/tree/main/percolation}} is laid out as a Palomar submission: a
statement file \texttt{Challenge.lean} importing only Mathlib, which defines the lattice, the measure, $\theta$, $\pc$ and the
proposition \texttt{PercolationContinuity d} and states two theorems with placeholder proofs; its twin \texttt{Solution.lean}, proving
the same theorems from the library \texttt{Percolation}; the comparator configuration \texttt{comparator.json}; and \texttt{AUDIT.md},
which lists what was mechanically checked and how to reproduce it.  To decide \emph{whether} the theorem is proved: read \texttt{Challenge.lean} (\cref{s:results}, box), run the build and comparator
as in \texttt{README.md}, inspect the axiom listing.  To see \emph{how}: read \cref{s:sketch} with the docstrings it names and \cref{s:guide}.
Where this text and a Lean statement differ, the Lean statement is what has been proved.

%=====================================================================================================
\section{Main results}\label{s:results}

\emph{Notation.}  A \emph{finite weighted graph} is a finite set $V$ with a weight $w_e\in[0,1]$ on every unordered pair $e$ of distinct
vertices (weight $0$ = absent edge); $\PP=\PP_w$ makes each pair $e$ open with probability $w_e$ independently (in Lean, $V=\mathrm{Fin}\,n$,
\lean{prodBernoulli w}).  $C_u$ is the vertex set of the open cluster of $u$, $\eC_u$ the set of open pairs inside it (the \emph{edge
cluster}); $\{u\con v\}=\{v\in C_u\}$ (\lean{openConn u v}), $\{u\con A\}=\bigcup_{a\in A}\{u\con a\}$; ``increasing'' means monotone for
inclusion, and increasing functions of the edge cluster $\eC_x$ include those of the vertex cluster $C_x=\{x\}\cup V(\eC_x)$;
$\E[\xi;B]:=\E[\xi\one_B]$ for an event $B$.

\begin{theorem}[$\theta(\pc)=0$ in all dimensions]\label{thm:main}
For every integer $d\ge2$, nearest-neighbour Bernoulli bond percolation on $\Zd$ satisfies $\theta(\pc)=0$, where
$\theta(p)=\PP_p(|C_0|=\infty)$ and $\pc=\inf\bigl(\{p\in[0,1]:\theta(p)>0\}\cup\{1\}\bigr)$.
\emph{Lean:} \lean{Percolation.Continuity.CSH.percolationContinuity_allDimensions} (for all $d$, $2\le d\to{}$\lean{PercolationContinuity d})
in \lean{Percolation/Continuity/MainTheorem.lean}, where \lean{PercolationContinuity d := theta (zdGraph d) 0 (criticalProbI d) = 0}
(\lean{Percolation/Literature/CriticalContinuity.lean}).  Comparator-checked form over Mathlib-only definitions:
\lean{BondPercolation.percolation_continuity} in \lean{Solution.lean}, against \lean{Challenge.lean}.  The case $d=3$, $\theta(\pc(\ZZ^3))=0$, is
\lean{CSH.percolationContinuity_three}, comparator-checked as \lean{BondPercolation.percolation_continuity_Z3}.\label{cor:Z3}
\end{theorem}
The `$\cup\{1\}$' only fixes the degenerate case $\theta\equiv0$; that this $\pc$ equals $\sup\{p:\theta(p)=0\}$, and the equivalence with
continuity of $\theta$ on $[0,1]$, are standard and \emph{not} part of the formal statement.

\begin{theorem}[additive gluing]\label{thm:AG}
Let $G$ be a finite weighted graph, $A$ a set of vertices, $o,b$ vertices and $t\ge0$.  If $\PP(a\con b)\ge1-t$ for every $a\in A$, then
\begin{equation}\label{eq:AG}
 \PP(o\con b)\ \ge\ \PP(o\con A)-t;\qquad\text{equivalently, for }A\neq\emptyset,\quad
 \PP(o\ncon b)\le\PP(o\ncon A)+\max_{a\in A}\PP(a\ncon b).
\end{equation}
\emph{Lean:} \lean{CSH.additiveGluing_holds : Statements.AdditiveGluing} (\lean{MainTheorem.lean}; statement in
\lean{Continuity/Statements.lean}).
\end{theorem}
For $|A|=1$ this is the union bound; the content is that it does not deteriorate with $|A|$.

\begin{corollary}[Kozma--Nitzan's Conjecture~3]\label{cor:conj3}
For every $\varepsilon>0$ there is $\delta>0$ \textup($\delta=\varepsilon/2$\textup) such that for every finite weighted graph, every $A$ and
all $o,b$: $\PP(o\con A)>1-\delta$ and $\PP(a\con b)>1-\delta$ for all $a\in A$ imply $\PP(o\con b)>1-\varepsilon$.
\emph{Lean:} \lean{CSH.kozmaNitzan_conjecture3_holds : KozmaNitzan2024_conjecture3} (\lean{MainTheorem.lean}); the statement, typed from
\citep[p.~15]{KozmaNitzan2024} in \lean{Literature/KozmaNitzanReduction.lean}, coincides with \lean{Statements.NearOneGluing} by
\lean{Iff.rfl} (\lean{nearOneGluing_iff_conjecture3}, \lean{Continuity/OfGluing.lean}).
\end{corollary}

\emph{The hierarchy: what the next definition encodes.}  For an increasing $f$ and any vertex $u$ the Harris inequality gives
$\Cov(f(\eC_x),\one\{u\in C_x\})\ge0$: nonnegative \emph{slack} between $f$ and ``the cluster of the \emph{owner} $x$ reaches $u$''.  The
hierarchy compares the slack seen at two \emph{observers}: at level~$0$ it says that the slack at $o$ is at least the fraction $\tau$ of the
slack at $v$, $\tau$ being the conditional probability that $o$ hangs on the cluster of $v$.  When \cref{ss:reduction} removes relays one at a
time, each removed vertex survives as a \emph{decoy} $z_j$ which has already explained the part $c_j(u)\times(\text{slack at }z_j)$ of the slack
at $u$; the level forms $\slf^j$ below subtract these parts, decoy by decoy, before the comparison is made.  The induction proving the hierarchy
(\cref{ss:csh}) runs in the complement of the cluster of an \emph{avoided set} $Y$, which is why the covariances are conditioned on
$\{x\ncon Y\}$ (the constants $c_j,\tau$ are nevertheless computed under $\PP$, \cref{rem:global}).  We write $\CSH(Y;x;Z;o,v)$ for the
assertion of \cref{thm:CSH} for the datum with avoided set $Y$, owner $x$, decoy list $Z$ and observers $o,v$ (in Lean \lean{CSHHolds w x Y D o v}).

\begin{definition}[datum, constants, level forms]\label{def:datum}
Let $V$ be finite with weights in $(0,1)$ on \emph{all} pairs (in Lean, on every element of \lean{Sym2 V}; diagonal pairs never matter).
A \emph{datum} is an \emph{owner} $x\in V$, an \emph{avoided set} $Y\subseteq V\setminus\{x\}$, a list $Z=(z_1,\dots,z_\ell)$ of distinct
\emph{decoys} ($\ell$ is the \emph{level}) and two distinct \emph{observers} $o,v$, the named vertices $x,o,v,z_1,\dots,z_\ell$ pairwise distinct
and outside $Y$.  Put $\nu:=\PP(\,\cdot\mid x\ncon Y)$, $Y_j:=\{x\}\cup Y\cup\{z_1,\dots,z_{j-1}\}$ ($1\le j\le \ell+1$), and define the
\emph{decoy constants} and the \emph{observer constant}
\begin{equation}\label{eq:constants}
 c_j(u):=\PP\bigl(u\in C_{z_j}\bigm|z_j\ncon Y_j\bigr)\ \ (1\le j\le \ell),\qquad \tau:=\PP\bigl(o\in C_v\bigm|v\ncon Y_{\ell+1}\bigr),
\end{equation}
computed under $\PP$, not $\nu$.  For a real function $\varphi$ on $V$ the \emph{level forms} are $\slf^0[\varphi]:=\varphi$,
\begin{equation}\label{eq:sl}
 \slf^j[\varphi](u):=\slf^{j-1}[\varphi](u)-c_j(u)\,\slf^{j-1}[\varphi](z_j)\ \ (1\le j\le \ell),\qquad
 \Marg[\varphi]:=\slf^\ell[\varphi](o)-\tau\,\slf^\ell[\varphi](v).
\end{equation}
\emph{Lean:} \lean{CSH.avoidConst}, \lean{decoyList}, \lean{obsConst}, \lean{slForm}, \lean{cshMarg} in \lean{Continuity/CSH/Defs.lean}.
\end{definition}

\begin{theorem}[the conditioned slack hierarchy]\label{thm:CSH}
For a datum as above and an increasing real function $f$ of sets of pairs, let
$\kappa_f(u):=\Cov_\nu\bigl(f(\eC_x),\one\{u\in C_x\}\bigr)$, $u\in V$.  Then
\begin{equation}\label{eq:CSH}
 \Marg[\kappa_f]\ \ge\ 0 .
\end{equation}
At level $0$ this reads
\begin{equation}\label{eq:level0}
 \Cov_\nu\bigl(f(\eC_x),\one\{o\in C_x\}\bigr)\ \ge\ \PP\bigl(o\in C_v\bigm|v\ncon\{x\}\cup Y\bigr)\cdot\Cov_\nu\bigl(f(\eC_x),\one\{v\in C_x\}\bigr),
\end{equation}
and at level $1$ with $Y=\emptyset$ it reads $\kappa_f(o)-c(o)\kappa_f(z)\ge\tau\,[\kappa_f(v)-c(v)\kappa_f(z)]$ with
$c(u)=\PP(u\in C_z\mid z\ncon x)$, $\tau=\PP(o\in C_v\mid v\ncon\{x,z\})$.
\emph{Lean:} \lean{CSH.cshHolds}, concluding \lean{CSHHolds w x Y D o v} (\lean{MainTheorem.lean}; Lean's \lean{D} is the decoy list $Z$;
\lean{CSHHolds}, \lean{cshMargin} in \lean{CSH/Defs.lean}).  The formal margin is \eqref{eq:CSH} times $\PP(x\ncon Y)^2>0$, each conditional covariance being carried in the
polynomial form $\PP(x\ncon Y)\E[f(\eC_x);x\ncon Y,x\con u]-\E[f(\eC_x);x\ncon Y]\PP(x\ncon Y,x\con u)$ (\lean{CSH.covD}); \lean{CSH.cshAll}
repackages the hypotheses over \lean{Fin n}.
\end{theorem}
For $Y=\emptyset$ and $f=\one\{v\in\cdot\,\}$, \eqref{eq:level0} is the Harris inequality: with $q:=\PP(x\con v)$ one has
$\kappa_f(o)=\PP(o,v\in C_x)-q\,\PP(o\in C_x)$, $\kappa_f(v)=q(1-q)$ and $\tau\kappa_f(v)=q\,\PP(o\in C_v,\,v\ncon x)$ (the factor
$1-q=\PP(v\ncon x)$ cancels against the conditioning in $\tau$), so that by $\{o\con\{x,v\}\}=\{o\in C_x\}\sqcup\{o\in C_v,\,v\ncon x\}$ the
inequality $\kappa_f(o)\ge\tau\kappa_f(v)$ reads $\PP(o\con\{x,v\},\,x\con v)\ge q\,\PP(o\con\{x,v\})$.  In general \eqref{eq:level0} says that, given that the owner avoids $Y$, the
correlation of $f$ with ``the owner reaches $o$'' is at least the part of its correlation with ``the owner reaches $v$'' that is transported
from $v$ to $o$ at the price $\tau$.

\begin{table}[t]
\newsavebox{\auditbox}\begin{lrbox}{\auditbox}\begin{minipage}{0.955\textwidth}\footnotesize
\textbf{Statements to audit.}  The kernel certifies that \lean{Solution.lean} proves the two theorems of \lean{Challenge.lean}; it does
not certify that they say ``$\theta(\pc)=0$ on $\Zd$''.  \lean{Challenge.lean} (namespace \lean{BondPercolation}) contains eleven definitions, one
lemma and the two theorem statements, listed here in full (the definitions are verbatim copies of the library's; the comparator requires
\lean{Solution.lean} to use identical ones).  A reader should check:
\begin{enumerate}[leftmargin=1.6em,itemsep=0pt,topsep=2pt]
\item \lean{Site d} $=$ functions \lean{Fin d} $\to\ZZ$; \lean{zdGraph d := SimpleGraph.hasse (Site d)}, adjacency $=$ covering relation of
 the coordinatewise order, i.e.\ $x\sim y$ iff $y=x\pm e_i$.  \emph{Check:} this is the nearest-neighbour graph
 (\lean{Pi.covBy_iff_exists_right_eq}, \lean{Order.covBy_iff_add_one_eq}).
\item \lean{BondConfig V := Set (Sym2 V)} (a configuration is its set of open pairs) and
 \lean{bondPercolation G p := ProbabilityTheory.setBernoulli G.edgeSet p}: edges of $G$ open independently with probability $p$, non-edges
 a.s.\ closed.  \emph{Check:} this is $\PP_p$.
\item \lean{openGraph}\ $\omega$ \lean{:= SimpleGraph.fromEdgeSet}\ $\omega$; \lean{openCluster}\ $\omega$\ \lean{x} $=$ the set of $y$ with
 \lean{(openGraph}\ $\omega$\lean{).Reachable x y}; \lean{percolatesAt x} $=$ the set of $\omega$ with \lean{(openCluster}\ $\omega$\ \lean{x).Infinite}.
 \emph{Check:} $\{|C_0|=\infty\}$ as a set of vertices.
\item \lean{theta G x p := (bondPercolation G p).real (percolatesAt x)}, a real number.
\item \lean{criticalProb G x} $=$ \lean{sInf} of $\{p\in\RR:\exists h\colon p\in[0,1],\ 0<{}$\lean{theta G x <p,h>}$\}\cup\{1\}$; the lemma
 \lean{criticalProb_mem_Icc} ($0\le\pc\le1$); \lean{criticalProbI d}, the same number as a point of $[0,1]$.  \emph{Check:} the $\cup\{1\}$
 convention; that the infimum form is the intended $\pc$.
\item \lean{PercolationContinuity d := theta (zdGraph d) 0 (criticalProbI d) = 0}.  \emph{Check:} $\theta$ at the origin.
\item the theorems \lean{percolation_continuity (d) (hd : 2 <= d) : PercolationContinuity d} and \lean{percolation_continuity_Z3 : PercolationContinuity 3}
 (placeholder proofs in \lean{Challenge.lean}, proofs in \lean{Solution.lean}); these two names are what \lean{comparator.json} compares.
 \emph{Check:} the hypothesis $2\le d$ (for $d=1$ the statement is false).
\end{enumerate}
\texttt{AUDIT.md} records \texttt{\#print axioms} (\texttt{[propext, Classical.choice, Quot.sound]} for all seven main theorems), the
\texttt{sorry} count (2, in \lean{Challenge.lean} by design), the absence of \texttt{axiom} declarations, and the comparator run.
\end{minipage}\end{lrbox}\noindent\fbox{\usebox{\auditbox}}
\end{table}

%=====================================================================================================
\section{Proof sketch}\label{s:sketch}

\paragraph{Roadmap.}
The proof has four steps, and it is easiest to read them from the lattice down to the finite graph.
\emph{Step~1} (\cref{ss:thm6}) is Kozma and Nitzan's theorem: if the near-one gluing statement (Conjecture~3) holds on every finite weighted
graph, then $\theta(\pc)=0$ on $\Zd$ for every $d\ge2$.  The mechanism is a renormalisation at \emph{fixed} $p$: an infinite cluster at $p$
lets one grow, box by box and with conditional probability close to one at each step, an infinite cluster inside a two-dimensional slab of
$\Zd$, gluing being exactly what carries the cluster from one box into the next; but no slab percolates at $\pc$.  This step is re-proved in
the library and nothing in it is new.  \emph{Step~2} (\cref{ss:conj3}) observes that Conjecture~3 follows at once from the additive gluing
inequality \eqref{eq:AG}.  \emph{Step~3} (\cref{ss:reduction}) derives additive gluing from the conditioned slack hierarchy: gluing through
a set $A$ of relays is proved by removing the relays one at a time, and what makes the induction close is a \emph{transfer} inequality
comparing a covariance seen from one observer with the same covariance seen from another; the removed relays do not disappear but survive as
``decoys'' whose influence is discounted, and the transfer inequality with $\ell$ decoys is precisely level $\ell$ of the hierarchy.
\emph{Step~4} (\cref{ss:csh}), the heart of the work, proves the hierarchy by induction on the level: conditioning on the cluster of the
avoided set splits the margin into a ``horizontal'' part, controlled by the Harris inequality, a negative-correlation inequality of
van den Berg--H\"aggstr\"om--Kahn and a new two-source inequality, plus one term per decoy which is a margin of \emph{lower} level.  The
simplest instance of level zero is the Harris inequality itself.  We present Step~4 before Step~3, because Step~3 consumes its notation.
Schematically (each arrow is a theorem of the development; beside it, what it consumes):
\begin{equation}\label{eq:chain}
\fbox{\parbox{0.9\textwidth}{\small\raggedright
 \textbf{$\theta(\pc)=0$ on $\Zd$ for all $d\ge2$}\quad(\cref{thm:main})\\[1pt]
 \hspace*{1.2em}$\Big\Uparrow$\ \ \parbox[t]{0.82\textwidth}{\S\ref{ss:thm6}: Kozma--Nitzan Thm.~6, re-proved; consumes Harris--Kesten ($d=2$), uniqueness of the infinite cluster, Barsky--Grimmett--Newman}\\[3pt]
 \textbf{Kozma--Nitzan Conjecture~3} on finite weighted graphs\quad(\cref{cor:conj3})\\[1pt]
 \hspace*{1.2em}$\Big\Uparrow$\ \ \parbox[t]{0.82\textwidth}{\S\ref{ss:conj3}: take $t=\delta=\varepsilon/2$}\\[3pt]
 \textbf{additive gluing} (AG)\quad(\cref{thm:AG})\\[1pt]
 \hspace*{1.2em}$\Big\Uparrow$\ \ \parbox[t]{0.82\textwidth}{\S\ref{ss:reduction}: peeling of relays, (S5D)$\Rightarrow$(S5)$\Rightarrow$(GEN)$\Rightarrow$(AG-loc); consumes vdBHK Thm.~1.3 and closure in the weights}\\[3pt]
 \textbf{conditioned slack hierarchy} (CSH), all levels\quad(\cref{thm:CSH})\\[1pt]
 \hspace*{1.2em}$\Big\Uparrow$\ \ \parbox[t]{0.82\textwidth}{\S\ref{ss:csh}: induction on the level; consumes the Gibbs sampler of vdBHK \S2, Gladkov's tree-Harris inequality, Ahlswede--Daykin}\\[3pt]
 \textbf{level $0$}, whose simplest instance is the Harris inequality; consumes Harris and vdBHK Thm.~1.4}}
\end{equation}
The labels are those of the Lean docstrings, each defined where it is first displayed: (S5) is a \emph{surplus-transfer} inequality between two
observers, (S5D) the same with decoys, (GEN) a first-relay lower bound for a general increasing functional and (AG-loc) its indicator case, a
localised union bound (\cref{ss:reduction}); (T), (U), (H) and Hpart are the three steps and the horizontal term of \cref{ss:csh}.  Every
subsection follows the pattern \emph{Goal / Idea / The inequality / Why it holds / In Lean}.

%-----------------------------------------------------------------------------------------------------
\subsection{From Conjecture 3 to \texorpdfstring{$\theta(\pc)=0$}{theta(pc)=0}: the theorem of Kozma and Nitzan (Step 1)}\label{ss:thm6}
\emph{Goal.}  Conjecture~3 $\Rightarrow$ $\theta(\pc)=0$ on $\Zd$ for every $d\ge2$ \citep[Theorem~6]{KozmaNitzan2024}.
\emph{Idea.}  Suppose $\theta(p)>0$.  Then large boxes are joined to far-away targets with probability close to one, and gluing says that such
near-certain connections can be chained \emph{without losing probability at each link}, however many candidate link points there are.  Chaining
them along a planar grid of boxes builds, at the same $p$, an infinite cluster confined to a thick two-dimensional slab; applied at $p=\pc$ this
contradicts the fact that slabs, which sit inside a half-space, do not percolate at $\pc(\Zd)$.  For $d=2$ the conclusion is classical.
Nothing in this step is new; we follow \citep[\S4]{KozmaNitzan2024} and record where the formal route differs from print.

\emph{The statement} is the implication \lean{KozmaNitzan2024_thm6} (\lean{Literature/KozmaNitzanReduction.lean}), proved as
\lean{KozmaNitzan2024_thm6_holds} (\lean{KozmaNitzanTheorem6SlabCritical.lean}).  \emph{Why it holds.}
\emph{$d=2$:} Harris' $\theta(\tfrac12)=0$ and Kesten's $\pc(\ZZ^2)=\tfrac12$ (\lean{percolationContinuity_two}, \lean{CriticalContinuity.lean};
\lean{harris_theta_half_holds}, \lean{kesten_criticalProb_Z2_holds}); the library gets $\theta(\tfrac12)=0$ from RSW estimates
\citep{Russo1978,SeymourWelsh1978} as in \citep[Ch.~3]{BollobasRiordanPercolation2006} and $\pc\le\tfrac12$ by self-duality
\citep[Thm.~(11.11)]{GrimmettPercolation1999} with sharpness in the form of \citep{DuminilCopinTassionEM2016} (\lean{perc_sharpness_holds}).
\emph{$d\ge3$:} with $S_k=\{0\le x_1\le k\}$ the slab and $\HH=\{x_1\ge0\}$, (a)~\emph{if Conjecture~3 holds and $\theta(p)>0$ then some slab
percolates at the same $p$} (\lean{KozmaNitzan2024_slabPercolation_holds}, \lean{KozmaNitzanTheorem6OfSlab.lean}); (b)~\emph{no slab percolates at
$\pc(\Zd)$} (\lean{theta_slab_criticalProb_zd_eq_zero_holds}); (a) at $p=\pc$ contradicts (b).

\looseness=-1 \emph{Half (b):} $\theta_{S_k}(\pc)\le\theta_\HH(\pc)$ by monotonicity in the graph (\lean{theta_induce_mono_holds}), and $\theta_\HH(\pc(\Zd))=0$ is
\citep[Thm.~1.1]{BarskyGrimmettNewman1991} in the form \citep[Thm.~(7.35)]{GrimmettPercolation1999} (\lean{BarskyGrimmettNewman1991_holds},
\lean{HalfSpaceProofs.lean}), proved for all $d\ge3$ by the block construction of \citep[\S7.3, Lemmas~(7.36), (7.52)]{GrimmettPercolation1999}
(\lean{BGNd.lemma_7_36C}, \lean{theta_pos_tall}, \lean{theta_pos_flat}): if $\theta_\HH(\pc)>0$, bricks on $\partial\HH$ would be ``good'' with
probability above a universal threshold, hence (finitely many edges) also at some $p'<\pc$, forcing $\theta(p')>0$.  \emph{Deviation:}
\citep[p.~25]{KozmaNitzan2024} and \citep[(7.38)]{GrimmettPercolation1999} take (b) from the Aizenman--Grimmett strict inequality $\pc(S_k)>\pc$
\citep{AizenmanGrimmett1991}; the formal proof avoids it (a.s.\ no open cluster of $\HH$ is both infinite and of bounded height,
\lean{BGNd.measure_percolatesVia_inter_heightLE_eq_zero}).  The slab theorem of \citep{DuminilCopinSidoraviciusTassion2016} is formalised
(\lean{DuminilCopinSidoraviciusTassion2016_holds}) but off the main path.

\looseness=-1 \emph{Half (a):} fix $d\ge3$ and $p$ with $\theta(p)>0$.  Conjecture~3 enters only through the \emph{target lemma} \citep[Lemma~10]{KozmaNitzan2024}
(\lean{KozmaNitzan.targetLemma}, \lean{KozmaNitzanTargetLemma.lean}): for every $\varepsilon$ there is $\delta$ such that, in any finite weighted graph
containing a copy $\Lambda$ of a lattice box with sub-box $\Lambda_0$ and a ``target'' $\mathsf T\subseteq\Lambda$ reachable from every point near
$\Lambda_0$ through a scaled \emph{hittable geometry} \citep[p.~16]{KozmaNitzan2024} (a region and a face such that, at this $p$, a large box is joined to
the scaled face inside the scaled region with probability tending to one), $\PP(o\con\Lambda_0)>1-\delta$ implies $\PP(o\con\mathsf T)>1-\varepsilon$ for all $o\notin\Lambda$.  The
proof explores the cluster of $o$ inwards until a shell of $\Lambda_0$ is crossed by many branches, attaches \emph{seeds} (bounded open edge sets at selected contact points \citep[p.~19]{KozmaNitzan2024}), uses
\emph{local uniqueness} in mesoscopic boxes to funnel everything into one cluster, contracts the revealed configuration to a finite graph, and applies
Conjecture~3 there with $\mathsf T$ wired to a point.  Inputs: the square-root trick \citep[(11.14)]{GrimmettPercolation1999} (\lean{sqrt_trick_holds})
and local uniqueness \citep[Lemma~7]{KozmaNitzan2024}, which Kozma and Nitzan take from Cerf \citep{Cerf2015} while noting that uniqueness of the
infinite cluster suffices; the library takes that route, re-proving uniqueness \citep{AizenmanKestenNewmanCMP1987,BurtonKeane1989},
\citep[Thm.~(8.1)]{GrimmettPercolation1999} (\lean{Grimmett1999_numInfiniteClusters_le_one_holds}).  Then come the \emph{corridor lemma}
\citep[Lemmas~11--12]{KozmaNitzan2024} (\lean{KozmaNitzanCorridor.lean}) and the renormalisation \citep[pp.~25--31]{KozmaNitzan2024}
(\lean{KozmaNitzanScheme}, \lean{-Steps}, \lean{-Theorem6.lean}): an exploration of macroscopic sites of $\ZZ^2\hookrightarrow\Zd$ inside the region
$L_r=\ZZ^2\times[-5r,5r]^{d-2}$, each examined site being good with conditional probability $>1-\varepsilon$ given the past \citep[(33)]{KozmaNitzan2024}
(\lean{KozmaNitzan.KSch.fail_bound}).  \emph{Deviation:} where \citep[p.~25]{KozmaNitzan2024} conclude by domination of a supercritical site
process (``we skip the details''), the library runs an explicit Peierls contour estimate with $\varepsilon\le2^{-32}$ (\lean{DualContours.lean},
\lean{HistorySiteRenormalization.lean}), obtaining an infinite cluster in $L_r$, hence in the slab $S_{10r}$ (\lean{KozmaNitzan.exists_slab_theta_pos_of_three_le}).
Conjecture~3 is applied only to finite graphs built from finitely many boxes with edges contracted or deleted.

%-----------------------------------------------------------------------------------------------------
\subsection{From additive gluing to Conjecture 3 (Step 2)}\label{ss:conj3}
\emph{Goal.}  \cref{thm:AG} $\Rightarrow$ \cref{cor:conj3}.  \emph{Why.}  Put $\delta=\varepsilon/2$; under the hypotheses of Conjecture~3,
\eqref{eq:AG} with $t=\varepsilon/2$ gives $\PP(o\con b)\ge\PP(o\con A)-\varepsilon/2>1-\varepsilon$.  \emph{In Lean:} \lean{additiveGluingSuffices_proof}
(\lean{AdditiveGluing/Suffices.lean}), \lean{CSH.conjecture3_of_csh}.
\emph{Illustration} (not part of the development).  On the $4$-cycle $o\,a_1\,b\,a_2$ with all weights $\tfrac12$ and $A=\{a_1,a_2\}$, enumeration of
the $16$ configurations gives $\PP(o\con A)=\tfrac34$, $\PP(a_i\con b)=\tfrac9{16}$, $\PP(o\con b)=\tfrac7{16}$; \eqref{eq:AG} with $t=\tfrac7{16}$ reads
$\tfrac7{16}\ge\tfrac34-\tfrac7{16}=\tfrac5{16}$.

We now turn to the finite-graph inequalities themselves: first the hierarchy (Step~4), then the derivation of additive gluing from it (Step~3).

%-----------------------------------------------------------------------------------------------------
\subsection{The conditioned slack hierarchy (Step 4)}\label{ss:csh}
\emph{Goal.}  \cref{thm:CSH}: for every datum and every increasing $f$, $\Marg[\kappa_f]\ge0$.  Throughout \S\S\ref{ss:csh}--\ref{ss:reduction} all
weights lie in $(0,1)$, so every event ``$u$ is joined to no vertex of a given set not containing $u$'' has positive probability; the lattice symbols
$p,d,k,S_k$ of \S\S\ref{s:intro}--\ref{ss:thm6} do not occur.  The symbols used more than once are:
\par\smallskip\noindent{\small
\begin{tabular}{@{}>{\raggedright}p{0.31\textwidth}>{\raggedright}p{0.49\textwidth}>{\raggedright\arraybackslash}p{0.14\textwidth}@{}}\toprule
symbol & meaning & defined\\\midrule
$x$;\ $Y$;\ $\nu=\PP(\cdot\mid x\ncon Y)$ & owner; avoided set; the conditional measure & Def.~\ref{def:datum}\\
$z_1,\dots,z_\ell$;\ $c_j$;\ $Y_j$;\ $X$ & decoys (level $\ell$); decoy constants; $\{x\}\cup Y\cup\{z_1,\dots,z_{j-1}\}$; marker set $\{x,z_1,\dots,z_\ell\}$ & Def.~\ref{def:datum}, \eqref{eq:constants}, \eqref{eq:U1}\\
$o,v$;\ $\tau$ & observers; observer constant $\PP(o\in C_v\mid v\ncon X\cup Y)$ ($v$ avoids everything in play) & Def.~\ref{def:datum}, \eqref{eq:constants}\\
$\Marg[\varphi]$ & the level form ``discount the decoys, then value at $o$ minus $\tau\times$ value at $v$'' & \eqref{eq:sl}\\
$\kappa_f(u)=\Cov_\nu(f(\eC_x),\one\{u\in C_x\})$ & conditioned covariance of $f$ with ``the owner reaches $u$'' & Thm.~\ref{thm:CSH}\\
$U=V\setminus C_Y$;\ $\pi$;\ $\PP_U$;\ $h_u(U)$, $\tau(U)$ & the \emph{world} (vertices outside the cluster of $Y$); its law under $\nu$; percolation on the graph induced on $U$; world covariance with ``$u\con X$'', world price & (T), \eqref{eq:U}, (H)\\
$N\subseteq U$;\ $\langle\psi\rangle^U_N$ & source set; ``explore $C_N$ inside $U$, evaluate $\psi$ on the rest'' & \eqref{eq:A2H}\\\bottomrule
\end{tabular}}\par\smallskip

\emph{Idea} (see also the paragraph before \cref{def:datum}).  Harris gives nonnegative slack $\Cov(f(\eC_x),\one\{u\in C_x\})\ge0$ at every
vertex $u$; gluing needs more, namely that the slack seen at $o$ is at least the fraction $\tau$ of the slack seen at $v$, $\tau$ being the chance
that $o$ hangs on $v$'s cluster, because that is what lets a connection be passed from $v$ to $o$ (level~$0$).  A relay $z$ removed in Step~3 has
already explained the part $c(u)\times(\text{slack at }z)$ of the slack at $u$; level~$\ell$ says the transfer survives $\ell$ such discounts.  The
induction works in the complement of the cluster of a vertex set $Y$, whence $\nu$; the constants $c_j,\tau$ are nevertheless global
(\cref{rem:global}).

\emph{Level $0$} is \eqref{eq:level0}: $\kappa_f(o)\ge\tau\,\kappa_f(v)$ with $\tau=\PP(o\in C_v\mid v\ncon\{x\}\cup Y)$.  The conditioning in $\tau$
cannot be dropped: on the path $o\!-\!x\!-\!v$ with both weights $\tfrac12$, $Y=\emptyset$ and $f=\one\{v\in\cdot\,\}$, the two edges are independent, so
$\kappa_f(o)=0$ while $\PP(o\con v)\,\kappa_f(v)=\tfrac14\cdot\tfrac14>0$; the conditioned price $\PP(o\in C_v\mid v\ncon x)$ is $0$.
\emph{Illustration} (not part of the development).  On the triangle $\{x,o,v\}$ with all weights $\tfrac12$, $Y=\emptyset$, $f=\one\{v\in\cdot\,\}$,
enumeration of the $8$ configurations gives $\kappa_f(o)=\tfrac7{64}$, $\kappa_f(v)=\tfrac{15}{64}$, $\tau=\tfrac13$, so the level-$0$ margin is
$\tfrac7{64}-\tfrac13\cdot\tfrac{15}{64}=\tfrac1{32}>0$.

\looseness=-1 \emph{Level $1$ with $Y=\emptyset$: the worked example.}  One decoy $z$; the statement is
$\kappa_f(o)-c(o)\kappa_f(z)\ge\tau\,[\kappa_f(v)-c(v)\kappa_f(z)]$ with $c(u)=\PP(u\in C_z\mid z\ncon x)$, $\tau=\PP(o\in C_v\mid v\ncon\{x,z\})$ and plain
covariances ($\nu=\PP$).  The mechanism of the whole proof is already visible here.  Put $X:=\{x,z\}$, $\mathcal E:=\{z\ncon x\}$,
$\rho:=\PP(\cdot\mid\mathcal E)$ and, for a vertex set $K\not\ni x$, $\Phi(K):=\E\bigl[f(\eC_x)-f(\eC^{-K}_x)\bigr]$, where $\eC^{-K}_x\subseteq\eC_x$ is
the owner's cluster in $G\setminus K$ (pairs meeting $K$ deleted); thus $\Phi\ge0$ and $\Phi$ is increasing in $K$.  Since
$\{u\con X\}=\{u\in C_x\}\sqcup(\{u\in C_z\}\cap\mathcal E)$ and $\{z\in C_x\}$ is the complement of $\mathcal E$, one gets
$\kappa_f(u)-c(u)\kappa_f(z)=\Cov(f(\eC_x),\one\{u\con X\})-\Cov\bigl(f(\eC_x),(\one\{u\in C_z\}-c(u))\one_{\mathcal E}\bigr)$, and the last
covariance equals $\E[f(\eC_x)(\one\{u\in C_z\}-c(u))\one_{\mathcal E}]$ because its second argument has mean $0$ by the choice of $c$.
Condition on $C_z=K$: on $\mathcal E$ the set $K$ misses $x$, and given $C_z=K$ the configuration off $K$ is a fresh percolation on
$G\setminus K$ (the Markov property, \ref{K2} below), so $\E[f(\eC_x)\mid C_z=K]=\E f(\eC^{-K}_x)=\E f(\eC_x)-\Phi(K)$; the constant $\E f(\eC_x)$
integrates to $0$ once more, and centring at $c(u)=\E_\rho\one\{u\in C_z\}$ turns what is left into a covariance:
\begin{equation}\label{eq:U1}
 \kappa_f(u)-c(u)\kappa_f(z)\;=\;\Cov\bigl(f(\eC_x),\one\{u\con X\}\bigr)+\PP(z\ncon x)\,\Cov_\rho\bigl(\Phi(C_z),\one\{u\in C_z\}\bigr).
\end{equation}
Taking \eqref{eq:U1} at $u=o$ minus $\tau$ times \eqref{eq:U1} at $u=v$, the level-one margin equals
\begin{multline*}
 \bigl[\Cov(f(\eC_x),\one\{o\con X\})-\tau\Cov(f(\eC_x),\one\{v\con X\})\bigr]\\[-2pt]
 +\PP(z\ncon x)\bigl[\Cov_\rho(\Phi(C_z),\one\{o\in C_z\})-\tau\Cov_\rho(\Phi(C_z),\one\{v\in C_z\})\bigr].
\end{multline*}
The first bracket is nonnegative by the \emph{set} version \eqref{eq:K6} below of level~$0$, for the marker set $X$ (Harris and \ref{K2}(b)).
The second bracket is exactly the level-$0$ margin of the datum with owner $z$, avoided set $\{x\}$, no decoys and the same observers, at the
functional $\Phi$---its observer constant $\PP(o\in C_v\mid v\ncon\{z\}\cup\{x\})$ is again $\tau$---so it is nonnegative by level~$0$.  Thus
level~$1$ at $Y=\emptyset$ already consumes level~$0$ with the \emph{non-empty} avoided set $\{x\}$ at a functional $\Phi$ that is not an
indicator: the induction forces the hierarchy to be stated for all avoided sets and all increasing $f$ (the peeling of \cref{ss:reduction}
independently needs $Y=R'\neq\emptyset$).  Once $Y\neq\emptyset$ the covariances live under $\nu$ and the computation above must be done
conditionally on the cluster of $Y$: that is where the \emph{worlds} and the two-source inequality below come from; at $Y=\emptyset$ there is
a single world, $\tau(U)=\tau$, and \eqref{eq:Htw} is an identity.

\emph{General form.}  \Cref{def:datum} performs the $\ell$ discounts successively---the $j$-th decoy is conditioned to avoid the owner, $Y$ and the
earlier decoys, which is what $Y_j$ records---and \cref{thm:CSH} asserts $\Marg[\kappa_f]\ge0$.  Two algebraic facts are used:
$\Marg[\varphi]=\sum_u\lambda(u)\varphi(u)$ over $u\in\{o,v,z_1,\dots,z_\ell\}$ (\lean{CSH.cshMarg_eq_sum_single}), and, with $\Marg'$ the form of
$z_2,\dots,z_\ell$ (same $c_2,\dots,c_\ell,\tau$),
\begin{equation}\label{eq:cons}
 \Marg[\varphi]=\Marg'[\varphi]-\varphi(z_1)\,\Marg'[c_1]
\end{equation}
(\lean{CSH.cshMarg_cons}).  Since $f\mapsto\Marg[\kappa_f]$ is linear and vanishes on constants, \eqref{eq:CSH} is a finite family of polynomial inequalities in the weights.

\emph{Inputs.}  The proof of \cref{thm:CSH} uses four published inequalities, all re-proved in the library:
\begin{enumerate}[label=\textup{(K\arabic*)},leftmargin=2.6em,itemsep=0pt,topsep=2pt]
\item\label{K1} the Harris inequality \citep{HarrisPCPS1960}, \citep[Thm.~(2.4)]{GrimmettPercolation1999};
\item\label{K2} conditional association \citep[Thms.~1.3--1.5, Remark~1]{VandenbergHaggstromKahn2005}: for vertex sets $V_1,V_2$, given
 $\{V_1\ncon V_2\}$, (a)~two increasing functions of $\eC_{V_1}:=\bigcup_{s\in V_1}\eC_s$ are positively correlated, (b)~an increasing function
 of $\eC_{V_1}$ and one of $\eC_{V_2}$ are \emph{negatively} correlated; and the Markov property of cluster exploration
 \citep[Lemmas~2.3--2.4]{VandenbergHaggstromKahn2005} (off the explored cluster of $V_1$ the configuration is a fresh percolation);
\item\label{K3} Gladkov's decision-tree Harris--Kleitman inequality \citep[Thm.~3.2]{Gladkov2024} (classical case
 \citep{HarrisPCPS1960,Kleitman1966}): if a decision tree run on two independent configurations $(\omega_1,\omega_2)$ reveals a random set
 $\mathcal S$ of coordinates, then $(\PP\otimes\PP)(\omega_1\in\mathsf A,\ \omega_1\!\to_{\mathcal S}\omega_2\in\mathsf B)\ge\PP(\mathsf A)\PP(\mathsf B)$
 for increasing events $\mathsf A,\mathsf B$, $\omega_1\!\to_{\mathcal S}\omega_2$ being $\omega_1$ on $\mathcal S$ and $\omega_2$ elsewhere; for the
 tree exploring the cluster $C_N$ of a set $N$, with exploration $\sigma$-algebra $\mathcal F_N$, this is $\Cov(\E[\xi\mid\mathcal F_N],\eta)\ge0$
 for increasing $\xi,\eta$: Harris survives stopping along an exploration (\lean{DecisionTree.PrW_mul_PrW_le_Pr2W_treeHK},
 \lean{TreeHarris.treeHarris_real}; also the kernel form of Gladkov--Zimin \citep[Ch.~5]{Gladkov2025Thesis}, \citep{GladkovZimin2024HK},
 \lean{GladkovZiminKernel.lean});
\item\label{K4} the Ahlswede--Daykin four functions theorem \citep{AhlswedeDaykin1978} (Mathlib's \lean{four_functions_theorem_univ}): for a
 product measure $\sigma$ on $2^{U'}$, $\alpha(s)\beta(t)\le\gamma(s\cup t)\delta(s\cap t)$ for all $s,t$ implies
 $\sigma[\alpha]\sigma[\beta]\le\sigma[\gamma]\sigma[\delta]$.
\end{enumerate}
Russo's formula, RSW and the BK inequality (in the library for $d=2$) are not used on this path.

\emph{Why \cref{thm:CSH} holds: worlds, and three steps.}  Fix a datum and an increasing $f$.  Under $\nu$ explore the cluster $C_Y$ and let $U:=V\setminus C_Y$, the \emph{world}, with
law $\pi$; on $\{x\ncon Y\}$ it contains $x$.  By the Markov property \ref{K2}, given $C_Y$ the configuration on $U$ is a fresh percolation on the graph
induced on $U$ (pairs meeting $C_Y$ deleted); write $\PP_U,\Cov_U$ for it, and set a world functional of $U$ to $0$ when a vertex it names is not
in $U$.  Let $\kappa^U(u):=\Cov_U(f(\eC_x),\one\{u\in C_x\})$ and let $\mathcal W[f]:=\E_\pi[\Marg[\kappa^U]]$ be the
\emph{world part}: the margin computed inside each world, but with the \emph{global} constants $c_j,\tau$, then averaged.  The proof is an induction
on the level $\ell$ (skeleton \lean{CSH.cshHolds_of_unfold}, \lean{AdditiveGluing/CSHInduction.lean}) in three steps: (T) reduces the margin to
the world part; (U) splits the world part exactly into a horizontal term plus one lower-level margin per decoy; (H) shows the horizontal term is
nonnegative.

\looseness=-1 \textsc{(T) Telescoping.}  \emph{Goal:} if $\mathcal W[f']\ge0$ for every increasing $f'\ge0$, then \eqref{eq:CSH} holds for every increasing $f$.
\emph{Idea:} conditioning on the world loses the part of the covariance carried by the world itself; that part is again a covariance of the same
shape for a new increasing functional, so one can iterate, and the iteration converges because the world is resampled from scratch with positive
probability.  \emph{Why:} given the world, $f(\eC_x)$ and $\one\{u\in C_x\}$ are functions of the percolation on $U$, so the law of total covariance under
$\nu$ gives $\kappa_f(u)=\E_\pi[\kappa^U(u)]+\Cov_\nu\bigl(g(U),\PP_U(u\in C_x)\bigr)$ with $g(U):=\E_Uf(\eC_x)$, and two applications of the tower
property ($g(U)$ is a function of $C_Y$, $\one\{u\in C_x\}$ of $\eC_x$) rewrite the last term as $\Cov_\nu((\mathcal Af)(\eC_x),\one\{u\in C_x\})$ with
$(\mathcal Af)(\mathcal K):=\E_\nu[g(U)\mid\eC_x=\mathcal K]$.  The constants $c_j,\tau$ do not depend on $f$, so applying the linear form $\Marg$:
the margin of $f$ is $\mathcal W[f]$ plus the margin of $\mathcal Af$.  Here $\mathcal A$ is the transition operator of the two-block Gibbs sampler
alternately resampling $C_Y$ given $\eC_x$ and $\eC_x$ given $C_Y$; $\mathcal Af$ is again increasing (a larger $\eC_x$ leaves less room for $C_Y$,
hence a larger world, and $g$ is increasing in $U$), and $\mathcal A^nf\to$ const since the chain regenerates with probability $\ge\prod(1-w_e)>0$
over the pairs leaving $Y$ (the one use of $w_e<1$); iterating, the margin of $f$ is $\sum_{i<n}\mathcal W[\mathcal A^if]$ plus a term tending
to $0$.  This is the scheme of \citep[proof of Thm.~2.1]{VandenbergHaggstromKahn2005} with Harris inside a world replaced by $\mathcal W\ge0$.
\emph{In Lean:} \lean{CSH.cshMargin_nonneg_of_within} (\lean{CSH/LemmaT.lean}), from \lean{BHK2006_multiMarkerCov_nonneg_of_within}
(\lean{Literature/TwoClusterGibbsCovariance.lean}).

\looseness=-1 \textsc{(U) Unfolding.}  \emph{Goal:} an exact formula for $\mathcal W[f]$.  \emph{Idea:} inside a world, ``the owner reaches $u$'' differs from
``$u$ is joined to the marker set $X$'' exactly by the events ``$u$ hangs on a decoy that misses the owner''; the level form was built so that these
differences, decoy by decoy, reassemble into margins of the \emph{same} hierarchy with that decoy as the new owner.  \emph{The identity:} with
$\mathcal E_j:=\{z_j\ncon Y_j\}$ and $\rho_j:=\PP(\cdot\mid\mathcal E_j)$,
\begin{equation}\label{eq:U}
\begin{split}
 \mathcal W[f]\;&=\;\mathrm{Hpart}\;+\;\sum_{j=1}^{\ell}\frac{\PP(\mathcal E_j)}{\PP(x\ncon Y)}\,
   \Marg^{(j)}\Bigl[u\mapsto\Cov_{\rho_j}\bigl(\Phi(C_{z_j}),\one\{u\in C_{z_j}\}\bigr)\Bigr],\\
 \mathrm{Hpart}&:=\E_\pi\bigl[h_o(U)-\tau\,h_v(U)\bigr],\qquad h_u(U):=\Cov_U\bigl(f(\eC_x),\one\{u\con X\}\bigr),
\end{split}
\end{equation}
where $\Marg^{(j)}$ is the form of $z_{j+1},\dots,z_\ell$ (same constants) and, for a vertex set $K$,
\begin{equation}\label{eq:Phi}
 \Phi(K):=\E\Bigl[\one\{x\ncon Y\text{ in }G\setminus K\}\bigl(f(\eC_x^{-C_Y^{-K}})-f(\eC_x^{-K})\bigr)\Bigr]
\end{equation}
($\eC_x^{-K}$, $C_Y^{-K}$: clusters computed in $G\setminus K$; in the first term only $C_Y^{-K}$ is deleted, not $K$): the mean excess of the
owner's cluster after deleting only the cluster of $Y$ grown without $K$, over the owner's cluster after deleting $K$.  $\Phi\ge0$ because on
$\{x\ncon Y$ in $G\setminus K\}$ the cluster $\eC_x^{-K}$ does not meet $C_Y^{-K}$ and therefore lies inside the cluster of the first term; $\Phi$ is increasing
in $K$ because enlarging $K$ shrinks $C_Y^{-K}$ (so the first cluster grows) and shrinks $\eC_x^{-K}$ (\lean{CSH.phiFun_nonneg}, \lean{phiFun_mono}); at
$Y=\emptyset$ it is the $\Phi$ of the worked example.  \emph{The point:} the $j$-th summand is $\PP(\mathcal E_j)/\PP(x\ncon Y)\ge0$
times the margin of the datum (owner $z_j$, avoided set $Y_j$, decoys $z_{j+1},\dots,z_\ell$, same observers) at the functional $\Phi$---a datum of
\emph{lower level} whose constants are literally $c_{j+1},\dots,c_\ell,\tau$---so the decoy terms are nonnegative by induction, and $\mathcal W[f]\ge0$
reduces to $\mathrm{Hpart}\ge0$.  At level $\ell=1$ with $Y=\emptyset$ there is one world $U=V$, and \eqref{eq:U} is the decomposition of the worked example
obtained from \eqref{eq:U1}.  \emph{Why \eqref{eq:U} holds.}  Fix a world $U$ and work under $\PP_U$.  Iterating the disjoint decomposition
$\{u\con X_i\}=\{u\con X_{i-1}\}\sqcup(\{u\in C_{z_i}\}\cap\{z_i\ncon X_{i-1}\})$, $X_i=\{x,z_1,\dots,z_i\}$, along the recursion \eqref{eq:sl}, as in
the derivation of \eqref{eq:U1}, gives the pointwise identity
$\slf^\ell[\one\{\cdot\in C_x\}](u)=\one\{u\con X\}-\sum_j\one\{z_j\ncon X_{j-1}\}\,\slf^{(j)}[\one\{\cdot\in C_{z_j}\}-c_j](u)+\mathrm{const}$, where
$\slf^{(j)}$ is the form of $z_{j+1},\dots,z_\ell$ (\lean{CSH.slForm_jn}); taking $\Cov_U(f(\eC_x),\cdot\,)$ (linear, zero on constants) and
averaging over $\pi$, the first term yields Hpart.  For the $j$-th term, the $\pi$-average of a world covariance $\Cov_U(f(\eC_x),\psi)$ is
$\E_\nu[(f(\eC_x)-g(U))\,\psi^U]$, $\psi^U$ being $\psi$ read in the world of the full configuration (Markov property at $C_Y$,
\lean{CSH.markov_merge_Y}); a decoy swallowed by $C_Y$ is isolated in its world, so its term is constant there and integrates to $0$ against the
residual (\lean{residual_orthogonal}), while a decoy outside $C_Y$ has the same cluster in the world as in $G$ and $\{z_j\ncon X_{j-1}$ in $U\}$
becomes $\mathcal E_j$.  Conditioning finally on $C_{z_j}=K$ (Markov property at $C_{z_j}$; on $\mathcal E_j$, $K$ misses $x$ and $Y$), the residual
times $\one\{x\ncon Y\}$ has conditional mean $-\Phi(K)$, the term $g$ producing the first half of \eqref{eq:Phi} (\lean{CSH.sum_phiIntegrand_eq}) and $f$
the second; centring at $c_j(u')=\E_{\rho_j}\one\{u'\in C_{z_j}\}$ gives the covariance of \eqref{eq:U}, the factor $\PP(\mathcal E_j)/\PP(x\ncon Y)$
coming from $\nu$ and $\rho_j$ (\lean{decoy_world_term}).
\emph{In Lean:} \lean{CSH.within_unfold}, \lean{subT_comb_nonneg} (\lean{CSH/UnfoldMain.lean}); \lean{decoy_world_term} (\lean{CSH/UnfoldDecoy.lean}).

\textsc{(H) Horizontal part.}  \emph{Goal:} $\mathrm{Hpart}\ge0$.  \emph{Idea:} inside each world the required transfer from $v$ to $o$ holds with the
world's \emph{own} price $\tau(U)$; the price actually charged is the global $\tau$, and although $\tau(U)-\tau$ has no sign in a given world, on average
over worlds the discrepancy helps.  \emph{The inequalities:} with the world's own price $\tau(U):=\PP_U(o\in C_v\mid v\ncon X)$ write
$\mathrm{Hpart}=\E_\pi[h_o(U)-\tau(U)h_v(U)]+\E_\pi[(\tau(U)-\tau)h_v(U)]$.  The first integrand is
nonnegative in every world by the \emph{set four-point transfer} (stated under $\PP$, applied under each $\PP_U$)
\begin{equation}\label{eq:K6}
 \Cov\bigl(f(\eC_x),\one\{o\con X\}\bigr)\ \ge\ \PP(o\in C_v\mid v\ncon X)\cdot\Cov\bigl(f(\eC_x),\one\{v\con X\}\bigr)\qquad(x\in X;\ o,v\notin X),
\end{equation}
and the second expectation is nonnegative by the \emph{two-source inequality}
\begin{equation}\label{eq:Htw}
 \PP(v\ncon X\cup Y,\ o\con v)\cdot\E\bigl[h_v(V\setminus C_Y);\,x\ncon Y\bigr]\ \le\ \PP(v\ncon X\cup Y)\cdot\E\bigl[\tau(V\setminus C_Y)\,h_v(V\setminus C_Y);\,x\ncon Y\bigr].
\end{equation}
Here $\E_\pi[\psi(U)]=\E[\psi(V\setminus C_Y);x\ncon Y]/\PP(x\ncon Y)$ for every world functional $\psi$, and $\tau=\PP(o\in C_v\mid v\ncon X\cup Y)$ because
$Y_{\ell+1}=X\cup Y$; so \eqref{eq:Htw} divided by $\PP(x\ncon Y)\,\PP(v\ncon X\cup Y)$ is precisely $\tau\,\E_\pi[h_v(U)]\le\E_\pi[\tau(U)h_v(U)]$.
\emph{Why \eqref{eq:K6} holds:} writing $f$ for $f(\eC_x)$, $\bar f:=f-\E f$, and using $\{o\con X\cup\{v\}\}=\{o\con X\}\sqcup(\{o\in C_v\}\cap\{v\ncon X\})$,
\begin{multline*}
 \Cov(f,\one\{o\con X\})=\Cov(f,\one\{o\con X\cup\{v\}\})-\E\bigl[\bar f\,\one\{o\in C_v\};v\ncon X\bigr]\\[-2pt]
 \ge\ -\PP(o\in C_v\mid v\ncon X)\,\E[\bar f;v\ncon X]\ =\ \PP(o\in C_v\mid v\ncon X)\Cov(f,\one\{v\con X\}):
\end{multline*}
the first covariance on the right is $\ge0$ by \ref{K1} and is dropped; the inequality is \ref{K2}(b) given $\{v\ncon X\}$ ($f$ is increasing in
$\eC_X$, $\one\{o\in C_v\}$ in $\eC_v$); the last step is $-\E[\bar f;v\ncon X]=\E[\bar f;v\con X]$.  For $X=\{x\}$, $Y=\emptyset$, \eqref{eq:K6} \emph{is} level zero \eqref{eq:level0}---the covariance
form of Kozma--Nitzan's two-relay mechanism \citep[Thm.~1]{KozmaNitzan2024}.  The two-source inequality \eqref{eq:Htw} is where \ref{K3} and \ref{K4}
enter.  \emph{In Lean:} \lean{CSH.hpart_nonneg} (\lean{AdditiveGluing/CSHHtwBridge.lean}), \lean{hpart_nonneg_of_htw}
(\lean{CSHHpart.lean}); \eqref{eq:Htw} is \lean{CSH.htw_world}, from \lean{CovTau.p1H_univ} (\lean{CovTau/A2H.lean}).

\looseness=-1 \textsc{The two-source inequality.}  \emph{Goal:} \eqref{eq:Htw}.  \emph{Idea:} it is the diagonal case of an inequality with two independent
``source sets'' whose clusters are explored, of the type introduced by van den Berg--Kahn and van den Berg--H\"aggstr\"om--Kahn for connection
probabilities; the novelty is that one of the four quantities is a covariance, which is not monotone in the configuration, and what replaces
monotonicity is a one-source bound obtained from Harris' inequality stopped along the exploration.  \emph{The inequality:} with worlds indexed by their vertex set $U$ as above, put
$\mathrm E_X(U):=\PP_U(o\in C_v,v\ncon X)$, $\mathrm M_X(U):=\PP_U(v\ncon X)$ (so $\tau(U)=\mathrm E_X(U)/\mathrm M_X(U)$) and $h:=h_v$; for a \emph{source set}
$N\subseteq U$ and any world functional $\psi$ let $\langle\psi\rangle^U_N:=\E_U[\psi(U\setminus C_N)]$: explore the cluster of $N$ inside $U$ and
evaluate $\psi$ on what is left.  Then \eqref{eq:Htw} is the diagonal $N=N'=Y$, $U=V$ of
\begin{equation}\label{eq:A2H}
 \langle\mathrm E_X\rangle^U_N\,\langle h\rangle^U_{N'}\ \le\ \langle\mathrm M_X\rangle^U_{N\cup N'}\,\langle\tau h\rangle^U_{N\cap N'}\qquad(N,N'\subseteq U)
\end{equation}
(indeed $\langle\mathrm E_X\rangle^V_Y=\PP(o\in C_v,\,v\ncon X\cup Y)$ and $\langle\mathrm M_X\rangle^V_Y=\PP(v\ncon X\cup Y)$ by the Markov property, and
$\langle h\rangle^V_Y=\E[h(V\setminus C_Y);x\ncon Y]$ since $h$ vanishes on worlds not containing $x$).
\emph{Why:} the input about $h$ is the pair of \emph{one-source} bounds, valid in every sub-world,
\begin{equation}\label{eq:starH}
 \langle h\rangle^U_N\cdot\PP_U(v\ncon X)\ \le\ \PP_U(v\ncon X,\,v\ncon N)\cdot h(U),\qquad \langle h\rangle^U_N\le h(U):
\end{equation}
exploring the cluster of a source set and recomputing the covariance in what is left loses on average at least the fraction
$\PP_U(v\con N\mid v\ncon X)$.  Proof of \eqref{eq:starH}: total covariance along $\mathcal F_N$; the term $\Cov(\E[f(\eC_x)\mid\mathcal F_N],\one\{v\con X\cup N\})$
that remains is $\ge0$ by \ref{K3}, and \ref{K2}(b) for the sets $X$, $\{v\}$ supplies the factor; antitonicity of $N\mapsto\langle h\rangle^U_N$ follows
from the second bound.  \emph{Base case of \eqref{eq:A2H}}, $N\cap N'=\emptyset$: by \eqref{eq:starH},
$\langle h\rangle_{N'}\le\PP_U(v\ncon N'\mid v\ncon X)\,h(U)$, and \ref{K2}(a) for the cluster of $v$ given $\{v\ncon X\}$ gives
$\langle\mathrm E_X\rangle_N\PP_U(v\ncon N'\mid v\ncon X)\le\tau(U)\langle\mathrm M_X\rangle_{N\cup N'}$; multiply.  \emph{Induction} on $|U|$, the scheme of
van den Berg--Kahn \citep{VandenBergKahn2001} and \citep[proof of Thm.~1.1]{VandenbergHaggstromKahn2005}: for $N_0:=N\cap N'\ne\emptyset$, condition on
the \emph{open star} $\mathfrak S\subseteq U':=U\setminus N_0$ of $N_0$ (vertices joined to $N_0$ by an open edge), whose law $\sigma$ is a product measure on
$2^{U'}$ and for which $C_N=N_0\cup C'_{(N\setminus N_0)\cup\mathfrak S}$ with $C'$ computed in $U'$; then each bracket in \eqref{eq:A2H} is a
$\sigma$-average, and \ref{K4} is applied on $2^{U'}$ to $\alpha(s)=\langle\mathrm E_X\rangle^{U'}_{(N\setminus N_0)\cup s}$,
$\beta(s)=\langle h\rangle^{U'}_{(N'\setminus N_0)\cup s}$, $\gamma(s)=\langle\mathrm M_X\rangle^{U'}_{((N\cup N')\setminus N_0)\cup s}$,
$\delta(s)=\langle\tau h\rangle^{U'}_{s}$, the pointwise hypothesis being \eqref{eq:A2H} in the smaller world $U'$ together with antitonicity in the
source set.  \emph{In Lean:} \lean{CovTau.a2H}, generic induction \lean{CovTau.metaA2_of_star} (\lean{CovTau/MetaA2.lean}), product law of the star
\lean{CovTau/A2Push.lean}, one-source bounds \lean{CovTau/StarH.lean}, antitonicity \lean{CovTau/MetaA2Anti.lean}.

\textsc{Level zero in the formal development} (an aside on the Lean route).  The formal proof of level~$0$ does not go through (T)+(H),
although those steps are formalised for every level: the development treats $\ell=0$ separately
(\lean{CSH.cshMargin_nil_nonneg}, from \lean{CovTau.markerDominanceAvoid}) and invokes (U) only for $\ell\ge1$; \eqref{eq:level0}, denominator-free and
for a general avoided set, is proved in \lean{Continuity/HullPort/} from a singleton-marker one-source bound of type \eqref{eq:starH} obtained from
\ref{K3} and \ref{K2} (\lean{HullPort.CE_holds}), by an induction on the avoided set in which the weight $\vartheta$ of one pair at a time is deformed
(\lean{HullPort.bernstein_step}, \lean{taQ_nonneg_of_Pv}): the quantity to be shown nonnegative has the form $\mathcal Q=\mathcal I_1\mathcal J_2-\mathcal J_1\mathcal I_2$
with $\mathcal I_i,\mathcal J_i$ affine in $\vartheta$, and a Bernstein-type identity expresses $\mathcal Q(\vartheta)$ on $[0,1]$ through
$\mathcal Q(0),\mathcal Q(1)$ (the pair deleted or contracted, covered by the induction) and a cross term whose sign is \ref{K2}(a) for the cluster
of $v$---without \ref{K4}.

\begin{remark}[why global constants]\label{rem:global}
Because the constants are computed under $\PP$, the sub-system generated by $z_j$ in \eqref{eq:U} is conditioned on exactly the event
defining $c_j$ and the later constants are unchanged, so the decoy terms are lower levels of the \emph{same} hierarchy; centring at world
constants instead would require a two-source inequality with two different covariance functionals, which is not available; the global
centring produces exactly \eqref{eq:Htw}.
\end{remark}

We now have the hierarchy for every owner, avoided set, decoy list and pair of observers.  The remaining issue is to turn this statement about
\emph{one} cluster and \emph{two} observers into a statement about an observer and a whole set of relays.

%-----------------------------------------------------------------------------------------------------
\subsection{From the hierarchy to additive gluing (Step 3)}\label{ss:reduction}
\emph{Goal.}  \cref{thm:CSH} $\Rightarrow$ \cref{thm:AG}.  \emph{Idea.}  Order the relays $a\in A$ by how valuable they are (here: by $\PP(a\con b)$)
and charge the failure of $o$ to reach $b$ to the \emph{first} relay that $o$ holds; this localised union bound (AG-loc) gives additive gluing
immediately, and it is the indicator case of a statement (GEN) about an arbitrary increasing functional $F$ of clusters: on the event that $o$ holds
some relay, $F(C_o)$ is on average worth at least the mean worth of the first relay held.  (GEN) is proved by removing the most valuable relay;
the cost of doing so is controlled by a \emph{surplus-transfer} inequality (S5) between two observers, and (S5) in turn is proved by the same
removal, the removed relays becoming decoys---which is where the hierarchy is consumed, twice per step.

\emph{The inequalities.}  Fix an increasing $F\ge0$ on vertex sets, $m_a:=\E F(C_a)$.  For a relay set $R$ list its elements so that $m$ is
non-decreasing (a \emph{compatible rank}); on $\{u\con R\}$ the \emph{first relay} $\iota(u)$ is the relay of least rank in $C_u$, and the events
$P^u_a=\{\iota(u)=a\}=\{u\con a\}\cap\bigcap_{a'\prec a}\{u\ncon a'\}$ partition $\{u\con R\}$.  The \emph{surplus} of $u$ over its first relay is
\begin{equation}\label{eq:Sur}
 \Sur_u(R):=\E\bigl[F(C_u);\,u\con R\bigr]-\sum_{a\in R}\PP(P^u_a)\,m_a=\E\Bigl[\bigl(F(C_u)-\min_{a\in C_u\cap R}m_a\bigr)\one\{u\con R\}\Bigr]
\end{equation}
(\lean{CSH.surplus}); the second form shows independence of the rank, and $\Sur_u(\{x\})=\Cov(F(C_x),\one\{u\in C_x\})$.  The links of
\eqref{eq:chain} between CSH and \eqref{eq:AG} are
\begin{align}
 \PP(v\ncon R)\cdot\Sur_o(R)\ &\ge\ \PP(o\con v,\ v\ncon R)\cdot\Sur_v(R) &&(v\notin R);\tag{S5}\label{eq:S5}\\
 \E\bigl[F(C_o);\,o\con A\bigr]\ &\ge\ \textstyle\sum_{a\in A}\PP(P^o_a)\,m_a;\tag{GEN}\label{eq:GEN}\\
 \PP(o\con A,\ o\ncon b)\ &\le\ \textstyle\sum_{a\in A}\PP(P^o_a)\,\PP(a\ncon b) &&(m_a=\PP(a\con b)).\tag{AG-loc}\label{eq:AGloc}
\end{align}
(S5) says the surplus seen from $o$ is at least that seen from $v$, discounted by the probability that $o$ is glued to $v$ while $v$ misses the
relays.  (S5D)$[R;Z;o,v]$ (relay set; decoy list; observers) is $\Marg[u\mapsto\Sur_u(R)]\ge0$ for a list $Z$ of decoys with constants as in \eqref{eq:constants} ($z_j$ avoiding
$R\cup\{z_1,\dots,z_{j-1}\}$, $v$ avoiding $R\cup Z$; \lean{CSH.surplusMargin}); (S5) is its case $\ell=0$ times $\PP(v\ncon R)$, and for $R=\{x\}$ it is
$\CSH(\emptyset;x;Z;o,v)$ at $\widehat F(\eC):=F(V(\eC)\cup\{x\})$---so with one relay, (S5D) \emph{is} the hierarchy.

\looseness=-1 \emph{Why $\CSH\Rightarrow$ \textup{(S5D)} (peeling).}  Weights in $(0,1)$; induction on $|R|$ for all decoy lists at once.  Let $a_s$ be the top relay,
$R'=R\setminus\{a_s\}$, $\mathcal D_s=\{a_s\ncon R'\}$, $\nu_s=\PP(\cdot\mid\mathcal D_s)$, $c_s(u)=\nu_s(u\in C_{a_s})$, $\gamma_s=\Cov(F(C_{a_s}),\one\{a_s\con R'\})$.
The bookkeeping fact: the constants of $[R;Z;o,v]$ are those of $[R';(a_s,Z);o,v]$, in which $a_s$ has become the \emph{first decoy}.  Then:
(P)~$\{u\con R\}=\{u\con R'\}\sqcup(\{u\in C_{a_s}\}\cap\mathcal D_s)$ gives
$\Sur_u(R)=\Sur_u(R')+\PP(\mathcal D_s)\Cov_{\nu_s}(F(C_{a_s}),\one\{u\in C_{a_s}\})-\gamma_sc_s(u)$ (\lean{surplus_erase_add});
($\gamma$)~$\gamma_s\le\Sur_{a_s}(R')$, as on $\{a_s\con R'\}$ the first relay of $a_s$ has smaller mean---the only use of compatibility
(\lean{kappa_le_surplus}); (AC)~$\Marg[c_s]\ge0$, by $\CSH(R';a_s;Z;o,v)$ at $\Psi_{\mathrm{iso}}(\eC)=\one\{\eC\ne\emptyset\}$, for which
$\Cov_{\nu_s}(\Psi_{\mathrm{iso}},\one\{u\in C_{a_s}\})=c_s(u)\nu_s(\eC_{a_s}=\emptyset)$ (\lean{covD_psiIso}); (N)~by \eqref{eq:cons},
$\Marg[\Sur(R')]-\Sur_{a_s}(R')\Marg[c_s]=\Marg'[\Sur(R')]$, the form with one relay fewer and one decoy more.  The middle term of (P) being
$\PP(\mathcal D_s)$ times the margin of $\CSH(R';a_s;Z;o,v)[\widehat F]\ge0$,
\begin{equation}\label{eq:peel}
 \begin{aligned}\Marg[\Sur(R)]\ &\ge\ \Marg[\Sur(R')]-\gamma_s\Marg[c_s]\\ &\ge\ \Marg[\Sur(R')]-\Sur_{a_s}(R')\,\Marg[c_s]\ =\ \Marg'[\Sur(R')]\ \ge\ 0\end{aligned}
\end{equation}
by induction.  Each step uses \cref{thm:CSH} twice; (S5) for $|R|=s$ consumes levels $0,\dots,s-1$.  \emph{In Lean:}
\lean{CSH.surplusMargin_nonneg_of_csh} (\lean{CSH/Peel.lean}, \lean{PeelTools.lean}).

\emph{Why \textup{(S5)}$\Rightarrow$\textup{(GEN)}.}  By \eqref{eq:Sur}, (GEN) says $\Sur_o(A)\ge0$.  Weights in $[0,1]$; induction on $|A|$ using
only (S5), \ref{K1}, \ref{K2}(a).  With
$a_s,A',\mathcal D_s,\gamma_s$ as above ($A$ in place of $R$) and $F_s=F(C_{a_s})$, (P) gives $\Sur_o(A)=\Sur_o(A')-\Delta$,
$\Delta=\E[(m_{a_s}-F_s)\one\{o\in C_{a_s}\}\one_{\mathcal D_s}]$, and by \ref{K2}(a) for the cluster of $a_s$ given $\mathcal D_s$, then ($\gamma$),
\begin{equation}\label{eq:deficit}
 \PP(\mathcal D_s)\,\Delta\ \le\ \PP(o\in C_{a_s},\mathcal D_s)\bigl(m_{a_s}\PP(\mathcal D_s)-\E[F_s;\mathcal D_s]\bigr)=\PP(o\in C_{a_s},\mathcal D_s)\,\gamma_s\ \le\
 \PP(o\in C_{a_s},\mathcal D_s)\,\Sur_{a_s}(A'),
\end{equation}
while (S5) for $A'$ with $v=a_s$ reads $\PP(\mathcal D_s)\Sur_o(A')\ge\PP(o\in C_{a_s},\mathcal D_s)\Sur_{a_s}(A')$; subtracting,
$\PP(\mathcal D_s)\Sur_o(A)\ge0$, which is (GEN) for $A$ if $\PP(\mathcal D_s)>0$.  If $\PP(\mathcal D_s)=0$ (possible with weights in $[0,1]$) then
$\Delta=0$ and $\Sur_o(A)=\Sur_o(A')\ge0$ by the induction hypothesis; this is its only use.  \emph{In Lean:}
\lean{AGloc.gen_firstRank_of_surplusTransfer} (\lean{AdditiveGluing/GenOfSurplusTransfer.lean}; two relays directly in \lean{GenPair.lean}).

\emph{Why \textup{(GEN)}$\Rightarrow$\textup{(AG-loc)}$\Rightarrow$\textup{(AG)}.}  At $F=\one\{b\in\cdot\,\}$, $m_a=\PP(a\con b)$ and $\E[F(C_o);o\con A]=\PP(o\con A,o\con b)$;
as the $P^o_a$ partition $\{o\con A\}$, (GEN) gives $\PP(o\con A,o\ncon b)\le\sum_a\PP(P^o_a)(1-\PP(a\con b))$; under $\PP(a\ncon b)\le t$ this is
$\le t\PP(o\con A)\le t$, and $\PP(o\con A)-\PP(o\con b)\le\PP(o\con A,o\ncon b)$.  \emph{In Lean:} \lean{AGloc.agloc_firstRank_of_gen},
\lean{additiveGluing_of_agloc_firstRank} (\lean{AdditiveGluing/OfAGloc.lean}).

\emph{Boundary cases.}  Peeling gives (S5) for weights in $(0,1)$ and $o\ne v$ outside $R$; $o=v$ and $o\in R$ are elementary
(\lean{surplus_nonneg_of_mem}), and weights in $[0,1]$ (absent and contracted edges) follow by closure: by the min-form in \eqref{eq:Sur} both sides
of (S5) are rank-free polynomials in the weights (\lean{CSH.surplus_eq_minForm}), so (S5) passes from $(0,1)^E$ to $[0,1]^E$, $E$ the set of pairs
(\lean{weights_le_of_forall_pos_lt_one}, \lean{AGloc.exists_rank_compat}; \lean{AdditiveGluing/OfSurplusTransfer.lean}, \lean{SurplusClosure.lean}).

\begin{remark}[Kozma--Nitzan's Conjecture~1 for three relays]\label{rem:mult}
The library proves the following (\lean{CovTau.kn_conj1_three}, \lean{Percolation/Continuity/CovTau/OfTA.lean}): on a finite weighted graph,
for vertices $o,b,a_1,a_2,a_3$ with $a_3\ne a_1$, $a_3\ne a_2$ and $\PP(a_1\con b)\le\PP(a_2\con b)\le\PP(a_3\con b)$, and every real
$t\le\PP(a_1\con b)$, one has $t\cdot\PP\bigl(\{o\con a_1\}\cup\{o\con a_2\}\cup\{o\con a_3\}\bigr)\le\PP(o\con b)$.  Up to relabelling the
relays so that $\PP(a_i\con b)$ is non-decreasing, this is Conjecture~1 \eqref{eq:KN1} for $|A|=3$.  The general Conjecture~1, and
Conjectures~2 and~4, are not addressed in this development.
\end{remark}

%-----------------------------------------------------------------------------------------------------
\subsection{Assembly: proof of Theorem~\ref{thm:main} from the pieces}\label{ss:assembly}
(i)~\cref{thm:CSH} holds: strong induction on the level $\ell$, the step being (T)+(U)+(H) of \cref{ss:csh} with (H) fed by the two-source
inequality \eqref{eq:Htw} (\lean{CSH.cshHolds} $=$ \lean{cshHolds_of_unfold} with \lean{within_nonneg_of_hpart} and \lean{hpart_nonneg}).
(ii)~By \cref{ss:reduction}, peeling turns \cref{thm:CSH} into (S5D), hence (S5), hence (GEN), hence (AG-loc), hence \cref{thm:AG} for weights in
$(0,1)$, and closure gives all weights (\lean{CSH.additiveGluing_of_csh}, \lean{CSH/AdditiveGluingOfCSH.lean}).
(iii)~By \cref{ss:conj3}, \cref{thm:AG} with $t=\varepsilon/2$ is \cref{cor:conj3} (\lean{CSH.conjecture3_of_csh}).
(iv)~By \cref{ss:thm6}, Kozma--Nitzan's Theorem~6 applied to \cref{cor:conj3} gives $\theta(\pc)=0$ on $\Zd$ for every $d\ge2$
(\lean{CSH.percolationContinuity_of_csh} via \lean{KozmaNitzan2024_thm6_holds}); $d=3$ is \lean{percolationContinuity_three}.
(v)~\lean{Solution.lean} transports \lean{percolationContinuity_allDimensions} along \lean{bridge} (\lean{Iff.rfl}: the two
\lean{PercolationContinuity} propositions unfold to the same term), and the comparator checks that the statement proved is that of
\lean{Challenge.lean}.\qed

%=====================================================================================================
\section{Guide to the formalization}\label{s:guide}
\begin{table}[t]\scriptsize\centering
\renewcommand{\arraystretch}{1.08}
\begin{tabular}{@{}>{\raggedright}p{0.28\textwidth}>{\raggedright}p{0.30\textwidth}>{\raggedright\arraybackslash}p{0.38\textwidth}@{}}
\toprule
folder / file & contents & key declarations\\\midrule
\lean{Challenge.lean}, \lean{Solution.lean} & statement file (Mathlib only); proved twin & \lean{BondPercolation.percolation_continuity}, \lean{percolation_continuity_Z3}, \lean{bridge}\\
\lean{Continuity/MainTheorem}, \lean{Statements}, \lean{OfGluing} & main theorems; gluing statements & \lean{CSH.cshHolds}, \lean{additiveGluing_holds}, \lean{kozmaNitzan_conjecture3_holds}, \lean{percolationContinuity_allDimensions}\\
\lean{Continuity/CSH/} (12 files) & \cref{def:datum}, \cref{thm:CSH}; (T), $\Phi$, (U); peeling & \lean{CSHHolds}, \lean{cshMargin}, \lean{cshMargin_nonneg_of_within}, \lean{within_unfold}, \lean{phiFun_mono}, \lean{surplusMargin_nonneg_of_csh}\\
\lean{Continuity/AdditiveGluing/} (13) & induction skeleton; (H); (S5)$\Rightarrow$(GEN)$\Rightarrow$(AG); weight closure & \lean{cshHolds_of_unfold}, \lean{hpart_nonneg}, \lean{gen_firstRank_of_surplusTransfer}, \lean{additiveGluing_of_agloc_firstRank}, \lean{surplusTransfer_of_nondegenerate}\\
\lean{Continuity/CovTau/} (19), \lean{HullPort/} (12), \lean{LowerTail/} (4) & two-source \eqref{eq:A2H}, one-source \eqref{eq:starH}; level-zero chain; tree-Harris for real functions & \lean{metaA2_of_star}, \lean{a2H}, \lean{p1H_univ}, \lean{markerDominanceAvoid}, \lean{taQ_nonneg_of_Pv}, \lean{treeHarris_real}, \lean{kn_conj1_three}\\
\lean{Literature/Basic}, \lean{CriticalContinuity}, \lean{LatticeModels/} & the model; product measures; weight continuity & \lean{zdGraph}, \lean{bondPercolation}, \lean{theta}, \lean{criticalProb}, \lean{prodBernoulli}, \lean{PercolationContinuity}\\
\lean{Literature/KozmaNitzan*} (16) & \citep[\S\S2--4]{KozmaNitzan2024}: Conj.~3, Thm.~6, Thms.~1, 3, 4, 7, 8, Lemmas 6--12 & \lean{KozmaNitzan2024_conjecture3}, \lean{KozmaNitzan2024_thm6_holds}, \lean{targetLemma}, \lean{KozmaNitzan2024_slabPercolation_holds}\\
\lean{Literature/HalfSpace*}, \lean{Tall*}, \lean{Flat*}, \lean{Slab*}, \lean{Uniqueness*}, \lean{DualContours} & BGN via \citep[\S7.3]{GrimmettPercolation1999}; DST; uniqueness; contours & \lean{BarskyGrimmettNewman1991_holds}, \lean{DuminilCopinSidoraviciusTassion2016_holds}, \lean{Grimmett1999_numInfiniteClusters_le_one_holds}\\
\lean{Literature/Harris*}, \lean{Kesten*}, \lean{RSW*}, \lean{Russo*}, \lean{PlanarDuality}, \lean{SharpnessDCT*} & $d=2$ & \lean{harris_theta_half_holds}, \lean{kesten_criticalProb_Z2_holds}, \lean{rsw_half_holds}, \lean{russo_formula_holds}, \lean{perc_sharpness_holds}\\
\lean{Literature/ConditionalPositiveAssociation*}, \lean{TwoCluster*}, \lean{TwoSet*}, \lean{DecisionTree*}, \lean{TargetExploration*}, \lean{GladkovZiminKernel} & \citep{VandenbergHaggstromKahn2005} Thms.~1.1--1.5 and \S2.1; \citep[Thm.~3.2]{Gladkov2024}; \citep[Ch.~5]{Gladkov2025Thesis}; Harris & \lean{BHK2006_clusterConditionalPositiveAssociation_holds}, \lean{BHK2006_twoSetConditionalAssociation}, \lean{BHK2006.harris}, \lean{PrW_mul_PrW_le_Pr2W_treeHK}, \lean{ED_ED_le_ED_diag}\\
\bottomrule
\end{tabular}
\caption{Module map.  Paths relative to \texttt{Percolation/}; declarations live in \texttt{Percolation.Continuity} (sub-namespaces \texttt{CSH}, \texttt{AGloc}, \texttt{CovTau}, \texttt{HullPort}) or \texttt{Percolation.Literature}.}\label{tab:map}
\end{table}


\emph{Size and pins.}  251 Lean files, 97{,}574 lines (87{,}136 non-blank): about 12{,}200 lines in 63 files under
\lean{Percolation/Continuity/} (this work) and 85{,}000 in 183 files under \lean{Percolation/Literature/}, plus \lean{Percolation.lean},
\lean{Util/Linter.lean}, \lean{Challenge.lean}, \lean{Solution.lean}, \lean{scripts/Axioms.lean}.  Toolchain \texttt{leanprover/lean4:v4.32.0};
Mathlib commit \texttt{81a5d257c8e4} (its tag \texttt{v4.32.0}), pinned \citep{MouraUllrich2021,Mathlib2020}.  \Cref{tab:map} is the module map.


\emph{Re-proved versus new.}  Everything under \lean{Literature/} re-proves published results from Mathlib (docstring tags
\texttt{[cite: Key, locator]} resolve in \texttt{references.bib}; a literature result is a \lean{def AuthorYear_result : Prop} discharged by
\lean{theorem AuthorYear_result_holds}; nothing is assumed).  Re-proved in this way are the results in the lower half of \cref{tab:map} (Kozma--Nitzan's Theorems 1, 3, 4, 6, 7, 8 and Lemmas
6--12 among them).  Everything under
\lean{Continuity/} is this work (a \texttt{[cite:]} tag there points to a statement a lemma specialises, not to a source of its proof): new
are \cref{thm:CSH}, \cref{thm:AG}, \cref{cor:conj3} and the chain between them, including \eqref{eq:S5}, \eqref{eq:GEN}, \eqref{eq:AGloc},
\eqref{eq:Htw}, \eqref{eq:A2H}.

\looseness=-1 \emph{Trusted base and checks.}  All seven main theorems depend only on \texttt{propext}, \texttt{Classical.choice}, \texttt{Quot.sound}
(\texttt{scripts/Axioms.lean}); \texttt{sorry} occurs twice, in \lean{Challenge.lean} by design; no \texttt{axiom} declarations, no unsafe code.  \texttt{lake exe cache get \&\& lake build} builds the default targets (about 6 minutes, per \texttt{AUDIT.md}); the CI script
runs the Palomar acceptance check---the Lean comparator\footnote{\url{https://github.com/leanprover/comparator}} with \texttt{lean4export}, the
independent \texttt{nanoda} kernel and the \texttt{landrun} sandbox---on \texttt{comparator.json}.  \texttt{formalization.yaml} records scope,
divergences and review status for the Palomar registry\footnote{\url{https://palomar-registry.org}}.  The docstrings are machine-written working notes, a map rather than an exposition.

%=====================================================================================================
\section{Limitations and open directions}\label{s:limits}
\emph{Other conjectures of Kozma--Nitzan.}  Conjecture~1 \eqref{eq:KN1} is proved only for $|A|=3$ (\lean{CovTau.kn_conj1_three}, stated
with its hypotheses in \cref{rem:mult}); the general Conjecture~1, Conjectures~2 and~4, and Questions~5--9 of \citep[\S5]{KozmaNitzan2024} are
not addressed in this development.  \emph{Rates.}  Nothing quantitative is proved: no bound on $\theta(p)$ as $p\downarrow\pc$, on
the tail of $|C_0|$ at $\pc$, or on $\delta(\varepsilon)$ beyond $\varepsilon/2$.  \emph{Other models.}  The lattice theorem is proved only for
nearest-neighbour bond percolation on $\Zd$; site percolation, other lattices and long-range models would need the reduction re-run there;
nothing new is proved about slabs.  \emph{Exposition.}  A conventional written proof refereed by experts does not yet exist.

\paragraph{Acknowledgements.}
\looseness=-1 The reduction that makes this result accessible is due to Gady Kozma and Shahaf Nitzan.  The Lean formalization and the first drafts of
this document were produced by Claude (Anthropic) working under my direction; I am responsible for the final text and for any errors.
I thank Ralph Furman and Levent Alp\"oge for reading and comments.  The formalization stands on Lean~4, on Mathlib and its community, and
on the Lean comparator, \texttt{lean4export}, \texttt{nanoda} and the Palomar template.

\bibliographystyle{amsplain-norepeat}
\setlength{\bibsep}{0pt plus 0.3pt}
{\footnotesize\bibliography{references}}

\end{document}
